\scanpage{50}
% archive foliation: 52
Générateurs et relations pour les groupoïdes et groupes de Teichmüller (spéciaux et ordinaires, évt. toriques\ldots)

Je m'attends à trouver trois types de présentations à propriétés complémentaires --- mais toutes à signification arithmétique :

a) Présentation par ``\uline{pantalons}'' ou ``\uline{triades}''. Elle revient à travailler sur les pts à l'infini ``les plus \unsure{cubiques|critiques}'' sur les multiplicités modulaires, et semble (d'après l'exemple de $\mathcal T_{0,4}^{+!}$) commode surtout pour exprimer des groupes du type $T_{g,\nu}^{+!}$ (ou $T_{g;I}^{+}$).

b) Présentation par ``boîte'' complète de pièces de Légo (\unsure{voire|voir} des boîtes encore plus complètes) à l'aide des \inserted{\uline{groupes d'}}\uline{automorphismes} \inserted{\uline{\unsure{stricts}}} de ces pièces, et certains ``chemins'' de déformation particuliers allant de certaines pièces à d'autres. Exemples frappants : $T_{0,4}^+$ (\unsure{connu!|connus!}), ou
\[
T_{1,1}^+\simeq\mathrm{SL}(2,\mathbb Z)\quad\text{engendré par }\rho,\sigma\quad\text{avec }\rho^{\unsure{3}}=\sigma^2,
\]
\[
\rho^6\ (=\sigma^4)=1.
\]

c) Présentations \uline{par involutions}, qui sont dans $\mathcal T_{g\nu}^{-}$. Exemple typique $\unsure{\mathbb T|T}_{11}\simeq\mathrm{GL}(2,\mathbb Z)$ \deleted{\unsure{eng}} (\deleted{\unsure{Origine}} aux Groupes cartographiques triangulés) engendré par trois réflexions $r_1,r_2,r_3$ avec
\[
(r_1r_2)^2=(r_2r_3)^3,\quad(r_1r_2)^4\bigl(=(r_2r_3)^6\bigr)=1.
\]
